\newpage
% ============================================================
\section{Parte E — Estilo del examen real: soluciones}

\begin{solucion}{E1 · Examen de 2025, pregunta 1 (taxis autónomos)}
\textbf{(a)} \emph{Desempeño:} tiempo de espera de los pasajeros, viajes atendidos, ingresos, km en vacío,
seguridad. \emph{Estados:} posición y estado de cada taxi (libre/ocupado/cargando, batería), demanda por
zona y hora, tráfico, hora del día. \emph{Perceptos:} GPS de la flota, solicitudes de la app, sensores de
tráfico, histórico de demanda. \emph{Acciones:} enviar un taxi libre a una zona, asignarle un viaje, dejarlo
esperando, mandarlo a cargar. \emph{Propiedades:} parcialmente observable (la demanda futura y el tráfico no
se ven), multiagente (otros conductores, peatones, competidores), estocástico (demanda y tráfico),
secuencial (la ubicación de hoy condiciona los viajes de mañana), dinámico, continuo (se suele discretizar
la ciudad en zonas y el tiempo en intervalos).

\textbf{(b)} \emph{Jugadores:} las compañías. \emph{Estado:} posiciones de todas las flotas + demanda.
\emph{Acciones:} reubicar vehículos (casi simultáneas). \emph{Utilidad:} participación de mercado o ingreso.
No es estrictamente de suma cero (la demanda total varía y hay zonas para todos) y tiene azar (demanda,
tráfico). Minimax supone un rival óptimo y totalmente hostil: da una estrategia robusta pero demasiado
conservadora. Expectimax (o expectiminimax, con nodos de azar para la demanda y un modelo probabilístico
del rival) maximiza la utilidad esperada, que es lo racional si el rival no es malicioso y la demanda es
aleatoria. Recomendación: \textbf{expectimax/expectiminimax} con profundidad limitada y una función de
evaluación (demanda esperada cubierta); minimax solo si se quiere garantizar el peor caso.
\end{solucion}

\begin{solucion}{E1 · Examen de 2025, pregunta 2 (A*)}
$h^*$: $G=0$, $C=3$, $B=4$, $A=5$, $S=\min(1+5,4+4)=6$. Admisible: $5\le6$, $3\le5$, $2\le4$, $1\le3$.

\textbf{(a)}
\begin{center}\small
\begin{tabular}{clll}
\toprule
Paso & Extrae ($g$, $f$) & Frontera después & Comentario\\
\midrule
1 & S (0, 5) & A(1+3=4), B(4+2=6) & \\
2 & A (1, 4) & C(3+1=4), B(6) & $f$ bajó de 5 a 4: síntoma de inconsistencia\\
3 & C (3, 4) & B(6), G(6+0=6) & \\
4 & B (4, 6) & G(6) & C vía B: $g=5>3$, no mejora\\
5 & G (6, 6) & — & meta\\
\bottomrule
\end{tabular}
\end{center}
Orden de expansión: S, A, C, B, G (si el empate B/G se rompe por G, B ni se expande). Camino S–A–C–G.

\textbf{(b)} Sí: cuesta 6 $=h^*(S)$; la alternativa S–B–C–G cuesta 8. Además, A* con $h$ admisible es
óptimo en árbol (argumento del ancestro en la frontera, C3a).

\textbf{(c)} Inconsistente en $S\to A$: $h(S)=5>c(S,A)+h(A)=1+3=4$. Con inconsistencia, $f$ puede decrecer a
lo largo de un camino y A* en grafo puede cerrar un nodo con un $g$ subóptimo; si no se reabre, el error
se propaga (ver C3 y E3). Aquí no ocurre porque cada nodo se extrae ya con su $g^*$. Para preservar la
optimalidad: (i) permitir la \emph{reapertura} de nodos cerrados cuando aparece un $g$ menor; (ii) búsqueda
en árbol (sin conjunto de cerrados); (iii) \emph{pathmax}: $f(n')\leftarrow\max(f(n'),f(n))$; (iv) sustituir por una
heurística consistente (p.\,ej.\ $h(S)\le4$).
\end{solucion}

\begin{solucion}{E1 · Examen de 2025, pregunta 3 (Connect-4 estocástico)}
\textbf{(a)} \emph{Estado} $s=(M,j)$: matriz $4\times4$ con celdas en $\{\text{vacía},X,O\}$ (respetando gravedad) y
jugador en turno $j\in\{\text{MAX},\text{MIN}\}$. \emph{Acciones:} $\textsc{Acciones}(s)=\{c: \text{la columna } c \text{ no está llena}\}$.
\emph{Transición:} sea $A(s)$ el conjunto de columnas disponibles (supuesto: el desvío es uniforme sobre
columnas no llenas):
\[P(c'\mid s,c)=\rho\,[c'=c]+\frac{1-\rho}{|A(s)|},\qquad c'\in A(s),\]
y $\textsc{Resultado}(s,c')$ coloca la ficha de $j$ en la fila vacía más baja de $c'$ y cambia el turno.
\emph{Terminal:} cuatro en línea (fila, columna o diagonal) o tablero lleno. \emph{Utilidad:} $+1$ gana MAX, $-1$
gana MIN, $0$ empate.

\textbf{(b)} Cada decisión va seguida de un nodo de azar:
\[V(s)=\begin{cases}
\textsc{Utilidad}(s) & s \text{ terminal}\\
\max_{c\in A(s)}\ \sum_{c'}P(c'\mid s,c)\,V(\textsc{Resultado}(s,c')) & \text{turno de MAX}\\
\min_{c\in A(s)}\ \sum_{c'}P(c'\mid s,c)\,V(\textsc{Resultado}(s,c')) & \text{turno de MIN}
\end{cases}\]
Con profundidad límite $d$ se sustituye $\textsc{Utilidad}$ por $\textsc{Eval}(s)$ al llegar a $d$ (p.\,ej.\ la
función lineal $3X_2+X_1-(3O_2+O_1)$ adaptada a líneas de 4). Si $\rho=1$ se recupera minimax. Costo
$O(b^m n^m)$ con $n$ resultados por nodo de azar. Como se promedian valores, la \emph{escala} de
$\textsc{Eval}$ importa (solo es invariante ante transformaciones afines positivas).
\end{solucion}

\begin{solucion}{E1 · Examen de 2025, pregunta 4 (CNF y resolución)}
\textbf{(a)} 1: todo estudiante que toma IA conoce A* ($\p{Estudiante}(x)\wedge\p{TomaIA}(x)\imp\p{Conoce}(x,\p{astar})$).
2: A* es un algoritmo. 3: A* es óptimo. 4: todo lo óptimo usa una heurística admisible. 5: todo lo
que usa una heurística admisible es correcto. 6: Bob es estudiante. 7: Bob toma IA.

\textbf{(b)} (1)+(6) $\{x/\p{bob}\}$: $\neg\p{TomaIA}(\p{bob})\vee\p{Conoce}(\p{bob},\p{astar})$; +(7): $\p{Conoce}(\p{bob},\p{astar})$.
(Por refutación: añadir $\neg\p{Conoce}(\p{bob},\p{astar})$ y un paso más da $\square$.)

\textbf{(c)} Sí. Refutación: añadir $\neg\p{Correcto}(\p{astar})$; con (5) $\{a/\p{astar}\}$: $\neg\p{UsaAdmisible}(\p{astar})$;
con (4) $\{a/\p{astar}\}$: $\neg\p{Optimo}(\p{astar})$; con (3): $\square$. La cláusula 2 no se usa.
\end{solucion}

\begin{solucion}{E2 · Bicicletas compartidas}
\textbf{(a)} Desempeño: viajes no atendidos (estación vacía o llena), km de camión, costo. Estados: número
de bicis y anclajes libres por estación, posición y carga de cada camión, hora. Perceptos: sensores de
anclaje, GPS de camiones, histórico de demanda. Acciones: mover un camión, cargar/descargar $k$ bicis.
Propiedades: parcialmente observable (demanda futura), multiagente (usuarios, otros camiones), estocástico,
secuencial, dinámico, discreto en bicis y continuo en tiempo/espacio (se discretiza). Es un problema de
\emph{optimización}: la ruta de cada camión puede plantearse como búsqueda (A* con costo en km), la
asignación de bicis por estación como CSP con restricciones de capacidad y, para instancias grandes,
búsqueda local o metaheurísticas (recocido, tabú) sobre planes completos.

\textbf{(b)} Jugadores: las dos empresas; turnos: decisiones periódicas (casi simultáneas; se modelan
alternadas); acciones: reubicar vehículos; utilidad: viajes capturados. No es de suma cero (la demanda
total cambia con la oferta). Expectiminimax con nodos de azar para la demanda si se quiere robustez frente
al rival; expectimax con un modelo probabilístico del rival si no es hostil. Árbol demasiado grande:
profundidad limitada + función de evaluación (demanda esperada cubierta), poda alfa-beta en la parte
determinista, profundización iterativa.
\end{solucion}

\begin{solucion}{E3 · A* con reapertura}
\textbf{(a)} $h^*$: $G=0$, $D=2$, $C=3$, $A=4$, $B=6$, $S=\min(2+4,1+6)=6$. Admisible. Violaciones:
$S\to B$: $5>1+1$; $A\to C$: $4>1+1$.

\textbf{(b)} Sin reapertura: S(0,5); B(1,2); C(4,5) \emph{se cierra con $g=4$} ($g^*=3$); A(2,6) genera C con $g=3$,
pero C está cerrado: se descarta; D(5,7); G(7,7). Devuelve \textbf{7}; el óptimo es S–A–C–D–G $=6$.

\textbf{(c)} Con reapertura: tras extraer A, C vuelve a la frontera con $(3,4)$; C(3,4); D(4,6); G(6,6):
devuelve \textbf{6}. El costo es reexpandir C (y lo que cuelga de él).

\textbf{(d)} Las dos violaciones involucran pares disjuntos ($S,B$ y $A,C$), así que hay que cambiar al menos
\textbf{dos} valores. Por ejemplo $h(B)=4$ y $h(C)=3$: todos los arcos cumplen ($5\le2+4$, $5\le1+4$, $4\le1+3$,
$4\le3+3$, $3\le1+2$, $2\le2$) y $h\le h^*$.
\end{solucion}

\begin{solucion}{E4 · Gato con viento}
\textbf{(a)} Estado: tablero $3\times3$ y jugador en turno. Acciones: casillas libres $L(s)$. Transición:
$P(s_{c'}\mid s,c)=\rho$ si $c'=c$ y $\frac{1-\rho}{|L(s)|-1}$ para cada $c'\in L(s)\setminus\{c\}$ (si $|L(s)|=1$, la marca
cae ahí con probabilidad 1), donde $s_{c'}$ es $s$ con la marca del jugador en $c'$ y el turno cambiado.
Terminal: tres en línea o tablero lleno; utilidad $+1/-1/0$.

\textbf{(b)} $V_d(s)=\textsc{Utilidad}(s)$ si es terminal; $\textsc{Eval}(s)$ si $d=0$; en otro caso
$\max_c$ (o $\min_c$) de $\sum_{c'}P(s_{c'}\mid s,c)\,V_{d-1}(s_{c'})$.

\textbf{(c)} $V(\textit{Atacar})=\rho(1)+(1-\rho)(-1)=2\rho-1$. $V(\textit{Bloquear})=\rho\cdot0+(1-\rho)\big(\tfrac12\cdot0+\tfrac12(-1)\big)=-\tfrac{1-\rho}{2}$.
Atacar conviene si $2\rho-1>-\tfrac{1-\rho}{2}\iff 4\rho-2>\rho-1\iff\rho>\tfrac13$. Minimax (azar como adversario):
ambas jugadas valen $-1$ en el peor caso; no distingue, mientras que expectimax sí usa $\rho$.
\end{solucion}

\begin{solucion}{E5 · KB de robots}
\textbf{(a)} 1: todo robot con batería se mueve. 2: todo lo que se mueve y está en ruta entrega. 3–4: $r_1$ es
un robot con batería. 5: todo dron vuela. 6: todo lo que vuela evita el tráfico. 7: $r_2$ está en ruta.

\textbf{(b)} (1)+(3) $\{x/r_1\}$: $\neg\p{Bateria}(r_1)\vee\p{Mueve}(r_1)$; +(4): $\p{Mueve}(r_1)$.

\textbf{(c)} $\p{Entrega}(r_1)$ \textbf{no} se sigue: falta $\p{EnRuta}(r_1)$ (solo $r_2$ está en ruta). $\p{EvitaTrafico}(r_1)$
\textbf{tampoco}: falta $\p{Dron}(r_1)$ o $\p{Vuela}(r_1)$. Modelo: dominio $\{r_1,r_2\}$ con $\p{Robot},\p{Bateria},\p{Mueve}$
verdaderos en $r_1$, $\p{EnRuta}$ solo en $r_2$ y $\p{Dron},\p{Vuela},\p{EvitaTrafico},\p{Entrega}$ falsos en todo: satisface
las 7 cláusulas y falsifica ambas conclusiones. (Por resolución, al añadir $\neg\p{Entrega}(r_1)$ solo se
obtiene $\neg\p{Mueve}(r_1)\vee\neg\p{EnRuta}(r_1)$ y luego $\neg\p{EnRuta}(r_1)$: nunca $\square$.)

\textbf{(d)} $\neg\p{Dron}(x)\vee\neg\p{BateriaBaja}(x)\vee\p{RegresaBase}(x)$ (sin Skolem).
$\forall x\,\p{Pedido}(x)\imp\exists y\,(\p{Robot}(y)\wedge\p{Entrega}(y,x))$ necesita Skolem $y=R(x)$:
$(\neg\p{Pedido}(x)\vee\p{Robot}(R(x)))\wedge(\neg\p{Pedido}(x)\vee\p{Entrega}(R(x),x))$.
\end{solucion}
